\documentclass{article}
\usepackage[utf8]{inputenc}
\usepackage[T1]{fontenc}
\usepackage{booktabs}
\usepackage{graphicx}
\usepackage{hyperref}
\usepackage{xcolor}
\definecolor{validated}{HTML}{047857}
\definecolor{refuted}{HTML}{B91C1C}
\definecolor{inconclusive}{HTML}{9A3412}
\title{Empirical Proof Record: VALIDATED}
\date{2026-05-17T17:32:44.332611+00:00}
\begin{document}
\maketitle
\section*{Empirical Proof Record --- \textcolor{validated}{\textbf{VALIDATED}}}
\textbf{Proof ID:} \texttt{4b824c5f2bb8f7b55e2b19cdd4450a9e\dots}\\
\textbf{Created:} 2026-05-17T17:32:44.332611+00:00
\subsection*{1. Claim}
\begin{quote}For Walker M4 events in a Sinew Phase 4-extended trajectory, the modulated-CW conservation ratio R\_mod EXCEEDS the seed-CW conservation ratio R\_seed by at least 0.05 — the substrate's Phase 148-182 SCAR-CW modulation chain encodes experience-accumulation at the cost of conservation. Both R values use the same (P + T + CW) three-term sum with std-normalized inputs; only the CW source differs (Phase 137 seed `\_p137\_scar\_w` vs the post-modulation `self.\_scar\_coherence\_weight`).\end{quote}
\begin{itemize}
\item \textbf{Operationalisation:} Read trajectory; compute R\_seed = median|Δ(P+T+CW\_seed)| / median|P+T+CW\_seed| at Walker M4 events. Compute R\_mod identically with CW = scar\_coherence\_weight\_modulated. Headline metric = R\_mod − R\_seed. Verdict on delta ≥ threshold. Secondary: per-event-class R\_seed vs R\_mod, cor(P+T, modulation\_delta) as counter-compensation signature, CW distribution statistics for both.
\item \textbf{Threshold:} \texttt{modulation\_disruption\_delta >= 0.05 dimensionless}
\item \textbf{H0:} modulation\_disruption\_delta < 0.05 — the substrate's modulation chain has been redesigned to preserve conservation; entry 991's finding regressed; the Phase 137 seed and the modulated value behave similarly at decision-points
\item \textbf{H1:} modulation\_disruption\_delta >= 0.05 — Phases 148-182 are counter-compensatory by design at decision-points; entry 991's finding holds; the substrate prioritizes memory formation (modulations push CW up when stress up) over strict conservation
\end{itemize}
\subsection*{2. Verdict}
\textcolor{validated}{\textbf{VALIDATED}} --- observed \texttt{0.07630155249669499}\\
Modulation-disruption measurement (500 cycles, 102 walker\_m4 events): R\_seed = 0.0873 [CI 0.0705, 0.1152], R\_mod = 0.1636 [CI 0.1519, 0.1991], delta = +0.0763. Counter-compensation signature cor(P+T, modulation\_delta) = +0.0405 (positive = modulations push CW up when stress up = anti-conservation by design).
\subsection*{3. Evidence}
\begin{tabular}{llrrl}\toprule
\textbf{Pillar} & \textbf{Statistic} & \textbf{Value} & \textbf{95\% CI} & \textbf{Library} \\\midrule
\texttt{sinew\_modulation\_disruption} & \texttt{modulation\_disruption\_delta} & \texttt{0.07630155249669499} & --- & ophamin.measuring.scenarios.sinew\_modulation\_disruption 0.9.0 \\
\bottomrule\end{tabular}
\subsection*{4. Pre-registration}
\begin{itemize}
\item Registered at: \texttt{2026-05-17T17:32:44.169158+00:00}
\item Config hash: \texttt{22f4999e3f8dad9f743ae5406d741096\dots}
\item Data hash: \texttt{fd2b8faaf10d26df1b01c6056b090ca8\dots}
\end{itemize}
\subsection*{5. Signature}
\texttt{eff73722f7476e11882497d530cc3eeccb6a5ee5bf9af39dfb31be6738d65a85}\\

\end{document}
